\chapter{Urinalysis: Ketones}
\section*{What Is This Test?}
Urine ketones testing detects ketone bodies (acetoacetate, beta-hydroxybutyrate, and acetone) in the urine. Ketone bodies are produced by hepatic mitochondrial fatty acid beta-oxidation and ketogenesis. They form when insulin is deficient or absent and counter-regulatory hormones (glucagon, cortisol, epinephrine, growth hormone) are high. This happens in: diabetic ketoacidosis (DKA), prolonged fasting/starvation, alcoholic ketoacidosis, very low-carbohydrate ketogenic diets, and pregnancy (starvation ketosis). In DKA, insulin deficiency causes uncontrolled lipolysis in adipose tissue. Free fatty acids are released, taken up by the liver, and converted to ketone bodies \cite{mcgarry1980regulation}. The buildup of acetoacetate and beta-hydroxybutyrate produces an increased anion gap metabolic acidosis. The urine ketone dipstick uses the Legal test (nitroprusside reaction). Sodium nitroprusside reacts with acetoacetate and acetone (but NOT beta-hydroxybutyrate) at an alkaline pH to produce a purple color. The dipstick detects acetoacetate with a sensitivity threshold of 5--10 mg/dL. In DKA, the main ketone body is beta-hydroxybutyrate (which is NOT detected by the dipstick). As DKA resolves with insulin therapy, beta-hydroxybutyrate is oxidized to acetoacetate. This paradoxically makes the urine ketone test stay positive or even rise (``ketone paradox'') despite clinical improvement. For this reason, direct measurement of serum beta-hydroxybutyrate is preferred for monitoring DKA treatment \cite{laffel1999ketone}.
\section*{Alternative Names}
\begin{itemize}
  \item Urine Ketones
  \item Ketonuria
  \item Ketone Bodies in Urine
  \item Acetoacetate Test
  \item Legal's Test
\end{itemize}
\section*{Principle}
Nitroprusside reaction (Legal test): Sodium nitroprusside + acetoacetate (or acetone) ---> purple-colored complex in alkaline pH. The dipstick does NOT react with beta-hydroxybutyrate. This is the main ketone body in DKA (accounting for 75--80\% of total ketone bodies in severe DKA). The detection threshold is 5--10 mg/dL. False-positive results occur with: highly pigmented urine, levodopa metabolites, sulfhydryl compounds including captopril and mesna (2-mercaptoethane sulfonate). False-negative results occur with improperly stored urine specimens in which acetoacetate has decarboxylated to acetone and evaporated.
\section*{History}
Wilhelm Petters first detected acetone in the urine of a diabetic patient in 1857. In 1883, Nicholas Legal described the sodium nitroprusside reaction for acetoacetate that still underlies the modern dipstick. In the early 1900s, Magnus-Levy established the relationship between ketone bodies and diabetic acidosis. After insulin was introduced in 1922, ketone measurement became essential for recognizing and monitoring diabetic ketoacidosis. Helen Free and Alfred Free developed the first practical urine ketone dipstick (Ketostix) in the 1950s. This made the nitroprusside test widely available at the bedside. It was recognized that the dipstick does not detect beta-hydroxybutyrate --- the main ketone in DKA --- and serum beta-hydroxybutyrate meters were developed in the 1990s for more accurate monitoring of ketosis.
\section*{Why This Test Is Ordered}
\begin{itemize}
  \item Diagnosing diabetic ketoacidosis in patients with hyperglycemia, and telling it apart from hyperosmolar hyperglycemic state \cite{kitabchi2001management}
  \item Monitoring DKA treatment, where a fall in serum beta-hydroxybutyrate confirms resolution even while urine ketones lag behind (ketone paradox) \cite{kitabchi2009hyperglycemic}
  \item Checking metabolic acidosis with an increased anion gap to identify a ketotic component
  \item Assessing ketosis in prolonged fasting, very low-carbohydrate (ketogenic) diets, and pregnancy (starvation ketosis)
  \item Diagnosing alcoholic ketoacidosis in patients with chronic alcohol use, vomiting, and starvation
  \item Monitoring patients on SGLT2 inhibitors who present with suspected euglycemic ketoacidosis
\end{itemize}
\section*{Primary vs Secondary Test}
Urine ketone testing is a primary point-of-care test for detecting ketosis in DKA. It is a rapid, inexpensive screen. Because it does not detect beta-hydroxybutyrate, serum beta-hydroxybutyrate is the preferred secondary (confirmatory and monitoring) test during treatment. The urine dipstick remains useful for screening. However, a negative urine ketone test does not exclude DKA when beta-hydroxybutyrate predominates.
\section*{How This Test Is Performed}
A freshly voided urine sample is collected in a clean container. The reagent strip is dipped into the specimen or passed through the midstream urine. The reagent pad is compared against the color chart after the manufacturer-specified time (typically 15--60 seconds). Results are reported semiquantitatively as negative, small, moderate, or large (approximate mg/dL of acetoacetate). For DKA monitoring, a capillary blood sample is tested with a beta-hydroxybutyrate meter \cite{tietz2023textbook}.
\section*{What Preparation Is Needed}
None. A fresh, well-mixed urine sample is preferred because acetoacetate is unstable. Standing at room temperature causes decarboxylation to acetone, which volatilizes and gives false-negative results. Testing should be done within 1--2 hours of collection.
\section*{Contraindications and Precautions}
\begin{itemize}
  \item No contraindications to urine collection. Precautions:
  \item (1) false-positive results occur with highly pigmented urine, levodopa metabolites, and sulfhydryl compounds (captopril, mesna);
  \item (2) false-negative results occur with old or improperly stored specimens;
  \item (3) highly acidic or strongly dilute urine can blunt the reaction;
  \item (4) the dipstick does not detect beta-hydroxybutyrate, so a negative urine test does not exclude DKA;
  \item (5) positive results should be read in the clinical context, because some drugs cross-react with the reagent pad.
\end{itemize}
\section*{Risks}
No physical risk; the test is noninvasive. The main concern is misinterpretation. This can happen when the specimen is stale or when the beta-hydroxybutyrate-predominant state of severe DKA is not recognized.
\section*{What Happens After the Test}
A positive urine ketone test in a diabetic patient triggers measurement of serum glucose, beta-hydroxybutyrate, electrolytes (anion gap, bicarbonate), and blood gas to confirm DKA. During treatment, serum beta-hydroxybutyrate is rechecked until it falls below 1.0--1.5 mmol/L and the anion gap normalizes. Urine ketones may stay positive for 12--24 hours after clinical resolution (ketone paradox). They should not delay the switch from intravenous insulin when other parameters have improved.
\section*{Understanding Results}
\subsection*{Negative}
No ketones detected. Normal metabolism without significant fatty acid oxidation or ketogenesis. Well-controlled diabetes.
\subsection*{Positive}
Small (5--15 mg/dL): Fasting, acute illness with decreased oral intake, vomiting, ketogenic diet, early DKA. Moderate (15--40 mg/dL): Early DKA, alcoholic ketoacidosis, pregnancy with inadequate caloric intake. Large (40--80--160+ mg/dL): DKA, severe alcoholic ketoacidosis. In DKA: positive urine ketones confirm the diagnosis. Serum beta-hydroxybutyrate is the preferred biomarker for measuring ketosis severity and monitoring treatment response (goal: beta-hydroxybutyrate <1.0 mmol/L before transition from IV insulin infusion). The urine ketone test may stay positive for 12--24 hours after resolution of DKA because of the ketone paradox.
\section*{Normal Results}
Negative (urine acetoacetate <5--10 mg/dL). Normal people do not excrete significant amounts of ketones. Pregnant women may have mild ketonuria with prolonged fasting or vomiting (hyperemesis gravidarum).
\section*{Follow-up Tests}
Serum beta-hydroxybutyrate (<0.6 mmol/L normal, >1.0 mmol/L = hyperketonemia, >3.0 mmol/L = severe ketosis/DKA), serum glucose, basic metabolic panel (anion gap, bicarbonate), serum osmolality, arterial or venous blood gas (pH, PCO2, HCO3), and serum lactate.
\section*{References}
\bibliographystyle{unsrtnat}
\bibliography{references}
\clearpage
\section*{Regulatory Status and Authorization Date}
This chapter describes a test category, not a specific commercial product. FDA, CDSCO, EU IVDR, and MHRA approval or registration dates therefore cannot be assigned without the assay manufacturer, product name, instrument, and intended use. For laboratory-developed tests, an individual agency approval date may not exist. See the Preface for the interpretation of regulatory status.

\clearpage
